\documentclass[a4paper]{article}
\usepackage{amsfonts}
\usepackage{amsmath}
\usepackage{amssymb}
\usepackage{graphicx}     

\begin{document}
  
\noindent \large Auteur: Anna Voronovska \\
\noindent \large Studierichting: Informatica\\
\noindent \large Oefeningenreeks 6B nr. 5.3\\

\medskip

\normalsize

Gegeven: $x_4 = -54, s_4 = -40$\\

Eerst $s_3$ zoeken:\\

$s_4 = s_3 + x_4$\\

$s_3 = s_4 - x_4 = -40 - (-54) = 14$\\

Als $x_4 = -54$ en $s_3 = 14$ dan:\\

(1) $\dfrac{s_3}{x_4} = \dfrac{14}{-54} = -\dfrac{7}{27} $\\

Formules:\\
 
$s_3 = x_1 \cdot \dfrac{q^3 - 1}{q - 1}$\\

$x_4 = x_1 \cdot q^3$\\

(2) $\dfrac{s_3}{x_4}
  = \dfrac{x_1 \cdot \dfrac{q^3-1}{q-1}}{x_1 \cdot q^3}
  = \dfrac{q^3 - 1}{q^3(q-1)}
  = \dfrac{(q-1)(q^2+q+1)}{q^3(q-1)}
  = \dfrac{q^2+q+1}{q^3}$\\

Nu (1) en (2) samen, krijgen wij:\\

$\dfrac{q^2+q+1}{q^3} = -\dfrac{7}{27}$\\

$27(q^2+q+1) = -7q^3$\\

$7q^3 + 27q^2 + 27q + 27 = 0$\\

Als $q = -3$:\\

$7(-3)^3 + 27(-3)^2 + 27(-3) + 27 = -189 + 243 - 81 + 27 = 0$\\

$q = -3$\\

Controle:\\

$x_1 = \dfrac{x_4}{q^3} = \dfrac{-54}{(-3)^3} = \dfrac{-54}{-27} = 2$\\

De rij: $2, -6, 18, -54, ...$\\

$s_4 = 2 + (-6) + 18 + (-54) = -40 $\\

Correct, dus $q=-3$


\begin{thebibliography}{999}
%\bibitem{a} Referentie
\end{thebibliography}
\end{document} 
